% %%%%%%%%%%%%%%%%%%%%%%%%%%%%%%%%%%%%%%%%%%%%%%%%%%%%%%%%%%%%
% CHAPTER 1 - INTRODUCTION (guideline 2.1)            ~7 pages
% Harvard citations with page numbers (guideline 4.2.1):
%   \parencite[1425]{krakowski2023}  -> (Krakowski et al., 2023, p. 1425)
%   \textcite[534]{gregory2021}      -> Gregory et al. (2021, p. 534)
% Statements are drawn from each article's abstract; the page given is
% the page on which the abstract appears (p. 1 for article-number journals).
% %%%%%%%%%%%%%%%%%%%%%%%%%%%%%%%%%%%%%%%%%%%%%%%%%%%%%%%%%%%%

\chapter{Introduction}
\label{ch:intro}

Few technologies have attracted as much attention from managers and researchers in recent years as artificial intelligence (AI). Online retailers use AI to personalise product recommendations, forecast demand and set prices, while digital platforms rely on it to match users, curate content and automate customer service. Yet the fact that a firm uses AI does not mean that it gains an edge over its competitors. As AI tools become widely available, the decisive question shifts from whether firms adopt AI to how they turn it into an advantage that rivals cannot easily copy, and whether this advantage lasts.

This thesis addresses this question from a dynamic capabilities perspective. Dynamic capabilities describe a firm's ability to sense new opportunities, seize them and transform its resource base as markets change \parencite[1319]{teece2007}. The thesis examines how AI capabilities, understood as the firm's ability to deploy and combine AI-specific resources, relate to these dynamic capabilities, and how both together contribute to sustainable competitive advantage in e-commerce and platform-based digital business. To answer this question, it systematically reviews the peer-reviewed literature on the topic. This chapter first outlines the background of the study, then defines the research problem and gap, states the research aim, questions and propositions, and finally explains the significance and structure of the thesis.

\section{Background of the Study}
\label{sec:background}

AI has moved from an experimental technology to a strategic resource that shapes how firms create and capture value. In the digital era, AI is no longer used only to automate isolated tasks: a systematic review of its strategic use identifies decision support, customer and employee engagement, automation, and new products and services as four sources of value creation \parencite[1]{borges2021}. Nevertheless, a further review of the business value of AI concludes that organisations still struggle to adopt and leverage AI and that a coherent understanding of how AI technologies create value is lacking \parencite[1709]{enholm2022}. For management research, AI therefore raises a classic strategic question in a new form: under which conditions does a technology become a source of superior and lasting performance?

This question is especially pressing in e-commerce and platform-based digital business. These business models generate large volumes of customer and transaction data and compete in digital markets whose dynamics depart from conventional competitive theory; here, the scale of a platform's user and complementor network and its distinct platform identity shape its competitive position \parencite[265]{cennamo2021}. AI adds a further layer to this competitive logic. \textcite[534]{gregory2021} argue that platforms can benefit from \enquote{data network effects}: the more a platform learns from the data of its users, the more valuable it becomes to each of them. AI capabilities thus have the potential to reinforce the scale-based advantages that already characterise digital platforms.

First empirical studies from e-commerce support this potential. In a survey of 394 e-commerce firms, \textcite[1]{chen2022} find that AI capability, conceptualised as a higher-order construct grounded in the resource-based view, affects firm performance indirectly through firm creativity, AI management and AI-driven decision making. More recently, \textcite[1]{bunea2026} show for electronic retailing that digital maturity drives AI implementation, which in turn generates benefits that are linked to customer experience and to differentiation- and cost-based competitive advantage. These findings suggest that AI can contribute to competitive outcomes, but they also indicate that the effect runs through organisational conditions rather than through the technology alone.

At the same time, AI is becoming more widely accessible. Cloud services and generative AI make advanced models available to almost any firm, which challenges the assumption that access to AI is in itself a scarce and valuable resource. Using a resource-based view, \textcite[1425]{krakowski2023} demonstrate that AI adoption can make traditional competitive capabilities obsolete while creating new sources of heterogeneity in how humans and machines work together. Similarly, \textcite[10]{li2025} show that generative AI capabilities contribute to competitive advantage mainly through service innovation, and that this effect depends on environmental dynamism. Notably, \textcite[10]{li2025} report the counterintuitive result that high environmental dynamism weakens some of these mediation effects, which suggests that the competitive value of AI capabilities is contingent on the market environment. Both studies point to the need to look beyond AI as a technology and to examine the capabilities that allow firms to deploy, combine and renew it.

Taken together, the background illustrates a tension that motivates this thesis. On the one hand, AI appears to be a powerful source of value in data-rich digital businesses. On the other hand, the growing availability of AI technologies and the dependence of their effects on complementary organisational factors call into question whether AI alone can be the basis of an advantage that competitors cannot easily imitate. Resolving this tension requires a theoretical lens that explains how firms build, deploy and continuously renew AI-related capabilities, which is why this thesis adopts a dynamic capabilities perspective.

\section{Research Problem and Research Gap}
\label{sec:problem}

Despite high expectations and substantial investments, many firms struggle to turn AI into measurable business value. \textcite[1]{alnofeli2026} note that although executives treat AI as a strategic priority, only a fraction of organisations report significant returns, particularly in customer-facing functions. Part of this difficulty lies in the nature of AI itself. Its probabilistic outputs, limited transparency and strong dependence on data create implementation challenges that require specific organisational capabilities \parencite[1549]{weber2023}. Based on 25 expert interviews, these authors derive four such capabilities: AI project planning and co-development to cope with the inscrutability of AI, and data management and AI model lifecycle management to cope with its data dependency. Such findings indicate that the value of AI depends less on the technology itself than on the organisational capabilities that surround it. The research problem of this thesis therefore concerns the gap between the adoption of AI and the creation of a competitive advantage that firms can sustain over time.

The resource-based view offers a first explanation for this gap. \textcite[1]{mikalef2021capability} identify the AI-specific resources that jointly create an AI capability and show that firms with a stronger AI capability achieve higher organisational creativity and performance. However, a resource-based explanation tends to be static. In markets characterised by rapid technological change, such as e-commerce, a lasting advantage depends on the ability to continuously sense opportunities, seize them and reconfigure the resource base \parencite[1319]{teece2007}. Dynamic capabilities theory therefore appears particularly suitable for explaining how AI-related resources are converted into sustainable competitive advantage.

A growing number of studies applies this perspective to AI. Based on case studies of large Norwegian organisations, \textcite[80]{mikalef2021b2b} identify micro-foundations through which AI supports sensing, seizing and transforming in business-to-business (B2B) marketing. Quantitative studies report that organisational resilience and strategic change \parencite[1]{liang2026} or strategic resilience \parencite[1]{wang2026} mediate the relationship between AI and firm outcomes. \textcite[1]{wu2026} add that identical AI investments can lead either to fleeting efficiency gains or to sustained advantages, depending on how managers interpret AI and which layer of dynamic capabilities they activate.

These studies also show that the question of sustainability is far from settled. \textcite[1]{liang2026} explicitly address whether AI enables firms to balance short-term economic returns with long-term sustainable growth, and they find that this effect varies considerably between firms. In a multi-method study of 357 early-stage ventures, \textcite[1]{wang2026} find that strategic resilience fully accounts for the relationship between AI innovation capability and entrepreneurial success, implying that AI investments create value only when they are converted into adaptive capacity. For e-commerce and platform businesses, in which technologies, customer expectations and competitors change rapidly, it therefore remains unclear whether and how AI capabilities translate into an advantage that persists over time.

Three gaps follow from this literature. First, the evidence is fragmented across industries: much of it stems from manufacturing, B2B marketing, banking or start-ups, while studies specific to e-commerce and digital platforms remain comparatively rare \parencites[1]{chen2022}[1]{bunea2026}. Second, many studies measure firm performance or short-term efficiency rather than a competitive advantage that is sustained over time, although these outcomes are conceptually distinct \parencite[1]{wu2026}. Third, existing reviews address the business value of AI \parencite[1709]{enholm2022} or AI-driven business model innovation, a field described as fragmented by varied conceptual lenses \parencite[1]{jorzik2024}. To the best of the author's knowledge, no systematic review has yet synthesised how AI capabilities, dynamic capabilities and sustainable competitive advantage are linked in e-commerce and platform-based digital business. This thesis addresses this gap.

\section{Research Aim and Questions}
\label{sec:objectives}

The aim of this study is to systematically review existing academic literature to examine how artificial intelligence capabilities contribute to sustainable competitive advantage in e-commerce and platform-based digital businesses, with particular attention to the role of dynamic capabilities, including sensing, seizing and transforming.

To achieve this aim, the thesis addresses three research questions:
\begin{enumerate}[label=RQ\arabic*:, leftmargin=1.5cm]
  \item How are artificial intelligence capabilities associated with competitive advantage in e-commerce and platform-based digital businesses?
  \item How do dynamic capabilities, particularly sensing, seizing and transforming, relate to the relationship between artificial intelligence capabilities and competitive advantage?
  \item What factors enable or constrain the development of sustainable competitive advantage through artificial intelligence capabilities in e-commerce and platform-based digital businesses?
\end{enumerate}

RQ1 examines the direct link between AI capabilities and competitive advantage. RQ2 places dynamic capabilities, the theoretical lens of this thesis, at the centre of the analysis. RQ3 asks under which conditions an AI-based advantage can be sustained or why it may erode. The three questions build on each other: RQ1 establishes whether a relationship exists, RQ2 explains the mechanism behind it, and RQ3 identifies the conditions under which the resulting advantage persists. All three questions are answered by systematically reviewing and synthesising peer-reviewed studies (Chapter~\ref{ch:methodology}); no primary data are collected. Throughout the thesis, competitive advantage and sustainable competitive advantage are treated as distinct outcomes: a short-term improvement in performance is not automatically interpreted as a sustainable advantage.

\section{Research Propositions}
\label{sec:propositions}

Since this thesis is based on a systematic literature review rather than on primary data, it formulates research propositions instead of statistically testable hypotheses. The propositions organise the review and are assessed in Section~\ref{sec:proposition-evaluation} as supported, contradicted or inconclusive in light of the reviewed evidence.

\begin{enumerate}[label=P\arabic*:, leftmargin=1.2cm]
  \item Artificial intelligence capabilities are associated with competitive advantage in digital business contexts, although the strength and nature of this relationship may vary across firms and business environments.
  \item Dynamic capabilities, particularly sensing, seizing and transforming, may help explain how firms convert artificial intelligence capabilities into competitive advantage.
  \item The sustainability of AI-enabled competitive advantage may depend on organisational and contextual factors, including the ability to integrate AI into business processes and to adapt to changing competitive conditions.
\end{enumerate}

P1 builds on resource-based studies that link AI capability to firm performance \parencite[1]{mikalef2021capability} and on survey evidence from supply chain executives that the assimilation of AI improves firm performance, with organisational and customer agility acting as partial mediators \parencite[1]{fossowamba2022}. P2 draws on dynamic capabilities theory \parencite[1319]{teece2007} and on studies in which capability-based mechanisms, such as service innovation, mediate the effect of generative AI capabilities on competitive advantage \parencite[10]{li2025}. P3 reflects the argument that AI implementation requires specific organisational capabilities \parencite[1549]{weber2023} and the finding that the effect of AI on long-term outcomes differs across firm characteristics such as size, industry and life-cycle stage \parencite[1]{liang2026}.

The propositions are deliberately formulated with caution, since the reviewed studies differ in context, design and measurement. A proposition is considered supported if the majority of the included studies, and especially those of high methodological quality, report findings consistent with it. It is considered contradicted if the evidence predominantly points in the opposite direction, and inconclusive if the findings are mixed or the number of relevant studies is too small to draw a conclusion. Evidence from e-commerce and platform businesses is weighted separately from supporting evidence from other industries.

\section{Significance of the Study}
\label{sec:significance}

From a theoretical perspective, the thesis contributes to research at the intersection of strategic management and information systems. It integrates the resource-based view and dynamic capabilities theory to explain not only whether AI capabilities are associated with competitive advantage, but also through which capabilities and under which conditions this advantage can be sustained. By separating evidence on performance, competitive advantage and sustainable competitive advantage, the review also responds to the conceptual ambiguity noted in recent work \parencite[1]{wu2026}. Furthermore, the thesis connects the emerging argument that AI shifts the sources of competitive advantage, because it both substitutes and complements human capabilities \parencite[1425]{krakowski2023}, with the platform-specific logic of data network effects. In doing so, it examines whether the dynamic capabilities perspective, which was developed for open economies with rapid innovation, also explains how advantages are built and renewed in AI-intensive digital markets. The resulting evidence-informed conceptual model can serve as a starting point for future empirical studies that test the proposed relationships with primary data.

From a managerial perspective, the findings are relevant for managers of e-commerce firms and digital platforms who decide on AI investments. Evidence from electronic retailing shows that digital maturity and AI implementation are key priorities for firms seeking AI-driven advantage \parencite[1]{bunea2026}. A structured synthesis of the literature can help managers identify which complementary capabilities need to be developed so that AI investments do not remain isolated technology projects. This is particularly relevant because, as \textcite[1]{alnofeli2026} emphasise, substantial investments in AI-powered systems frequently fail to generate meaningful business value.

From a methodological perspective, the thesis applies a transparent and replicable review procedure. The search strategy, selection criteria, screening steps and quality assessment are documented in detail, so that other researchers can reproduce the review and extend it as new studies are published. Given the rapid growth of research on AI in business, such a structured overview of the current state of evidence is a useful basis for future empirical work.

The scope of the thesis is limited to peer-reviewed academic literature on AI capabilities, dynamic capabilities and competitive advantage. Technical questions of AI system design, legal and ethical regulation, and public-sector applications are not examined. Studies from other industries are included only as supporting evidence and are identified as such in the analysis.

\section{Structure of the Thesis}
\label{sec:structure}

The thesis is organised in five chapters. Following this introduction, Chapter~\ref{ch:literature} reviews the literature on AI capabilities in digital business, introduces the resource-based view and dynamic capabilities theory, and discusses the sensing, seizing and transforming capabilities. It then examines the link between AI capabilities and competitive advantage, the specific conditions of e-commerce and platform businesses, and concludes with the research gap and the conceptual framework of the study.

Chapter~\ref{ch:methodology} explains the research philosophy and design of the systematic literature review. It documents the literature search strategy, the inclusion and exclusion criteria, the procedure for study selection and data extraction, the quality assessment and evidence synthesis, and the methodological limitations.

Chapter~\ref{ch:results} presents the characteristics of the included studies and synthesises the evidence on AI capabilities and competitive advantage, on AI capabilities and dynamic capabilities, and on the mediating role of dynamic capabilities. On this basis, it evaluates the research propositions and discusses the findings in relation to the literature reviewed in Chapter~\ref{ch:literature}. Finally, Chapter~\ref{ch:conclusion} answers the research questions, outlines the theoretical contributions and managerial recommendations, and discusses limitations and avenues for future research.
